\section{Benchmark design}
\label{sec:method}

\subsection{Backbones}
\label{sec:backbones}

35 monomeric PDB entries, chain A, 56--415 residues, downloaded and validated
by \texttt{prep\_targets.py} (CA count 50--450, parseable chain A). The list is
fixed in \texttt{data/targets30.txt} so that future arms cannot change the
problem set.

The composition of the set --- length distribution, experimental method and
resolution --- is tabulated in Table~\ref{tab:backbones} in the appendix.

The set is deliberately broad in fold space but \textbf{skewed small} (more
than half are under 100 residues), which matches the monomer regime where
inverse folding is usually benchmarked but not the 150--300 residue range that
dominates real design campaigns. Five NMR entries are included because they are
standard small-monomer benchmarks; ProteinMPNN is known to recover native
sequences poorly on NMR structures, and those backbones contribute some of the
low-scTM cells in every arm. Both points are revisited in the limitations.

\subsection{Arms}
\label{sec:arms}

\begin{table}[htbp]
  \centering
  \small
  \begin{tabular}{@{}llp{2.2in}@{}}
    \toprule
    Family & Arm & Mechanism \\
    \midrule
    published & ProteinMPNN \texttt{v\_48\_020} & baseline (a seed-1 replicate is run only to measure repeatability, and is not an arm) \\
    published & SolubleMPNN \texttt{v\_48\_020} & retrained without transmembrane entries \\
    published & HyperMPNN \texttt{v48\_020\_epoch300\_hyper} & retrained on thermophile AlphaFold2 structures \\
    published & MembraneMPNN global-label, fed ``soluble'' & conditioned input \\
    published & LigandMPNN \texttt{v\_32\_020} on backbone atoms & extra atom-context encoder \\
    guided & surf $\lambda$10 & per-residue apolar penalty on exposed sites \\
    guided & agg $\lambda$10 & windowed (5-residue) exposed apolar penalty \\
    guided & charge $\lambda{+}3$ & constant per-AA charge tilt \\
    guided & surf10 / charge$-2$ & two-property tilt (acidic branch) \\
    guided & surf10 / charge$+2$ & two-property tilt (basic branch) \\
    \bottomrule
  \end{tabular}
  \caption{Benchmark arms and their mechanisms.}
  \label{tab:arms}
\end{table}

Sixteen designs per (arm, backbone), temperature 0.1, fixed seed. The main
table below is this 16-design run; a predecessor run at 8 designs per cell is
kept in \texttt{results/AGGREGATE.md} and is used only for the temperature
ablation.

\subsection{Guidance: the AO-ARM closed form}
\label{sec:guidance}

ProteinMPNN samples residues in a random order and conditions on any subset of
known positions, i.e.\ it \emph{is} an any-order autoregressive model.
ProteinGuide's update for AO-ARMs multiplies the per-position categorical by a
likelihood ratio. For a property that is linear in amino-acid composition,
$\phi(a)=\sum_a n_a w_a$, that ratio is position-independent:
\[
  \log \tilde{p}(x_i{=}a \mid x_S)
  = \log p_\theta(a \mid x_S) - \lambda \cdot w_a + \mathrm{const},
\]
so guidance is \textbf{exactly} a constant per-amino-acid logit bias --- no
predictor calls, no approximation. ProteinMPNN's sampler computes
\texttt{softmax(logits/T + bias/T)}, hence we store $b = \lambda T$ to obtain an
effective logit shift $\lambda$. Surface hydrophobicity uses a per-residue
version, $b_{i,a} = -\lambda \cdot h_a \cdot s_i$, where $s_i \in \{0,1\}$
marks exposure computed from the \textbf{input backbone} (C$\alpha$ neighbour
count, median split), never from the design. Multi-property arms add the biases.

\paragraph{Boundary.}
Non-linear or neighbour-dependent objectives (aggregation patches beyond the
windowed proxy, $\pi$-stacking, stability) need ProteinGuide's TAG approximation
and are out of scope here.

\subsection{Metrics}
\label{sec:metrics}

\begin{itemize}
\item \textbf{scTM}: TM-score of the ESMFold-refolded design~\cite{esmfold}
  against the input backbone, at valid sites.
\item \textbf{recovery}: identity to the native sequence.
\item \textbf{diversity}: mean pairwise Hamming within the design set of one
  (arm, backbone) cell (16 designs in the main run, 8 in the \S4.1b ablation).
\item \textbf{controllability slope $k$}: OLS slope of achieved property on
  requested $\lambda$.
\item \textbf{leakage}: $\sum_{j\neq\mathrm{target}} |\Delta\phi_j|$
  normalised by $|\Delta\phi_{\mathrm{target}}|$.
\end{itemize}

\subsection{Statistical protocol}
\label{sec:stats}

Unit of analysis is the \textbf{backbone}, never the design. For every arm
comparison we form per-backbone means, then a paired bootstrap over backbones
($B = 10{,}000$) and report the 95\% percentile interval. Repeatability is
estimated by re-running ProteinMPNN with a different sampling seed on all 35
backbones (same weights, same temperature), which bounds the resolution of
every comparison. Designs are never pooled across backbones. Nine comparisons are made
against one baseline, so we report raw bootstrap $p$-values \textbf{and}
Holm-adjusted $p$-values and treat an effect as confirmed only if it survives
Holm at $\alpha = 0.05$.

\textbf{Power.} The seed-1 replicate gives the paired per-backbone standard
deviation of a null comparison directly, sd = 0.0240 scTM over 35 backbones.
With that dispersion the minimum detectable effect at 80\% power is
\textbf{0.0114 scTM (1.35\% relative)} for a single uncorrected comparison and
\textbf{0.0147 scTM (1.74\%)} once the Holm threshold for nine comparisons is
applied. This number is the honest interpretation of every null result in this
paper: the data exclude effects larger than about 0.015 scTM, not effects of
zero. It also explains why the guided arms' point estimates ($+0.003$ to
$+0.009$) cannot be called improvements, and why HyperMPNN's $-0.064$, which is
four times the corrected threshold, is the only effect of a size this benchmark
can resolve.
